\chapter{Conclusions}
\label{chap:conclusion}

\section{Contributions and main findings}

This thesis derives an Expanded Welch--Satterthwaite test for comparing
categorical MI across independent populations. Its mathematical
contribution is the observation-level sensitivity of the estimated MI
variance, accounting for changes in cell probabilities, margins and
pointwise MI. The variance of this sensitivity gives a first-order
sampling-variance approximation, which is used with moment matching to
choose MI-specific degrees of freedom for a Student \(t\) distribution.

\subsection{RQ1: Accuracy of variance approximations}

The mean and sampling variance of \(\Vhat(P)\) need separate checks.
The predicted sampling variance can be close to its observed value while
the mean estimate remains above \(V(P)\). In the moment target in
Equation~\eqref{eq:design-moment-target}, the mean is squared in the
numerator. Replacing the inflated sampling mean by the smaller population
value therefore lowers the population first-order degrees of freedom,
even when the sampling-variance approximation is accurate. The mean
estimated component degrees of freedom can fall further below the moment
target when population probabilities are replaced by sample proportions.
The size of these discrepancies depends on the population and sample size.

Near independence, the first-order prediction \(\tau^2(P)/n_P\) for
\(\Var\{\Vhat(P)\}\) is small and quadratic sampling terms remain
substantial. The calculation in
Appendix~\ref{sec:independence-boundary} explains the observed scale of
the estimated component degrees of freedom. These results
explain why the first-order approximation becomes inadequate near
independence and why substituting sample probabilities requires its own
evaluation, even when the population approximation is accurate.

\subsection{RQ2.1: False-positive rates under equal MI}

Expanded Welch brings false-positive rates closer to 5\% in the selected
small-sample settings A and B, both at common MI 0.02 nats
(Table~\ref{tab:redesign-focus-settings}). Both tests return valid results
in every repetition at these settings, so the lower rejection rates
reflect the larger cutoff from the Student \(t\) distribution. Across the
broad grid, Expanded Welch also lowers rates that are already below 5\%.
Its larger cutoff helps when it corrects excessive rejection and increases
conservatism when rejection is already too infrequent.

A false-positive rate near 5\% does not establish that the estimated
degrees of freedom accurately match the denominator's sampling moments.
At A, the standard error understates the observed sampling variation of
the MI difference, while the low estimated degrees of freedom increase
the rejection cutoff. These errors can compensate in the false-positive
rate. At B, the average squared standard error exceeds the observed
variance of the MI difference and the difference is biased, so the same
explanation does not apply unchanged.

The sparse same-skew null illustrates how error in the standard error can
remain when the expected difference bias cancels. Identical population
tables at equal sample sizes have the same expected MI estimation bias.
However, the average \(\SEhat^2\) is almost twice the observed
\(\Var(\Deltahat)\) in the sparse copying case, and both tests are
conservative. Unequal sample sizes can also leave unequal residual biases
in otherwise identical populations. Reliable testing therefore requires
accurate estimation of both the MI difference and its uncertainty.

\subsection{RQ2.2: Nominal and calibrated power}

Expanded Welch has lower nominal power than Normal Wald at the selected
small-sample settings A and B. Because the two tests also have different
false-positive rates, the calibrated comparison adjusts their cutoffs to
match rejection at a specified equal-MI population before comparing power.
At common MI 0.02 nats, Expanded Welch gains power over Wald at A and B.
Moving calibration to the upper equal-MI population while holding the
alternative fixed reverses the gain at A and two null SDs, and at B and
four null SDs (Table~\ref{tab:redesign-calibration-sensitivity}). At C,
nominal power differences are small and calibrated differences remain
practically small under both calibrations.

The calibrated power advantage depends on the population used to fit the
cutoff. Matching false-positive rates at the lower common MI does not
ensure that the two tests have the same false-positive rate at the upper
common MI. Expanded Welch's sample-dependent degrees of freedom can change
the ordering of datasets by their p-values, but the observed gains do not
establish a stable power advantage across the equal-MI null.

Changing the cutoff also leaves errors in the estimated MI difference
unchanged. At B, positive estimation bias initially offsets the negative
true difference as \(I(Q)\) increases. The estimated difference moves
towards zero, so rejection initially falls even though the population MI
values move farther apart. This distortion remains after calibration.

The sparse \(16\times16\) different-skew comparison shows a more severe
failure. Its positive difference bias exceeds the true negative difference
at the 0.01-nat alternative, so almost every rejection gives the wrong
estimated ordering. Equation~\eqref{eq:redesign-alternative-bias-shift}
explains the sign reversal. The four-cell construction also reproduces
the sparse-table failure. A larger rejection cutoff cannot correct that bias.

\section{Limitations}

The derivation assumes independent categorical observations, fixed
categories, positive population cell probabilities and positive
first-order MI variance. Paired samples would require a covariance term
in the difference variance. Growing alphabets and continuous MI estimators
have different approximation errors. Exact independence requires the
second-order argument in Appendix~\ref{sec:independence-boundary}.

The findings concern the tested populations, association constructions and
alternative paths. The power settings were selected using earlier
false-positive results and evaluated on independent samples. Each curve
keeps \(I(P)=0.02\) and increases \(I(Q)\) along the copying family.
Both tested constructions use equal row and column margins within each
table and align association locations between populations. Other directions
of MI change, unequal row and column margins within a table and other
locations of dependence remain untested. The results therefore support
further investigation of Expanded Welch rather than its use as a general
replacement for Normal Wald.

Empirical calibration uses known simulation populations. The two calibration
checks do not bound rejection across the entire equal-MI null or supply
a calibration procedure for arbitrary observed tables. The pointwise
confidence intervals quantify simulation uncertainty for the specified
populations; the calibrated power intervals also include uncertainty from
fitting the cutoffs.

The online-shopping application illustrates how to report the estimated
MI difference and its direction, standard error, p-value and Expanded Welch
degrees of freedom. Its unknown population MI values prevent a direct
check of false-positive rates. The application therefore demonstrates
interpretation of the test outputs, while the simulation evidence
establishes their behaviour under known populations.

\section{Future work}

A first extension is to derive the second-order mean contribution to the
variance estimator and evaluate a corresponding degrees-of-freedom
correction. The observed variation across MI and margins indicates that
the correction should depend on the population table. A fixed coefficient
based only on table size would leave errors in some settings. Any
estimated correction should be evaluated through the resulting test's
false-positive rates and power.

Improved MI bias and standard-error estimators could address the sparse-table
failures. Their effects should be assessed separately by changing one
estimate while holding the others fixed, then measuring null rejection,
power and estimated direction. Near independence, the quadratic sampling
approximation may support a more accurate comparison of MI differences.
Any likelihood-ratio or resampling procedure must preserve the equal-MI
null even when the complete population tables differ.

Reversed power paths, stronger dependence at small sample sizes,
additional association constructions, unequal margins and unequal sample
sizes would assess whether power rankings persist across a wider range
of populations. Simulations from plausible application populations could
also examine sensitivity to rare cells and distinct equal-MI tables,
with conclusions tied to the populations chosen for that assessment.
